\section{问题分析}

\subsection{问题一分析}

附件数据包含267名孕妇共1082次纵向检测记录，每名孕妇平均约4次观测，构成非平衡面板结构。同一孕妇的重复测量之间存在显著相关性（组内相关系数ICC约0.60），若使用普通最小二乘回归将低估标准误并产生虚假显著性。Y染色体浓度随孕周的变化呈非线性特征，BMI的稀释效应可能与孕周存在交互作用。因此需要采用能同时处理个体聚类效应和非线性趋势的纵向数据模型。

\subsection{问题二分析}

检测时点的选择面临双重约束。过早检测时胎儿DNA浓度偏低，检测灵敏度不足；过晚检测则压缩了后续临床干预的时间窗口。不同BMI组别的孕妇DNA浓度曲线形态和达标时间存在系统性差异，需要在分组基础上构建同时反映准确性风险和延迟风险的目标函数，通过优化求解各组的最优时点。

\subsection{问题三分析}

引入年龄、体重、身高等多因素后，模型维度增加，变量间可能存在共线性（如体重与BMI高度相关）。达标率的评估需要在多维体征空间上进行条件概率计算。测量误差从自变量端传入模型后，会通过非线性映射放大或缩小，需利用Delta方法或蒙特卡洛模拟量化其对预测结果的影响。

\subsection{问题四分析}

女性胎儿缺乏Y染色体标记，无法直接复用前三问的浓度模型。染色体非整倍体的检测依赖对全基因组测序覆盖度的统计比较，Z值衡量目标染色体读数相对参考分布的偏离程度。正常与异常样本的Z值分布存在重叠区域，单一阈值难以兼顾灵敏度与特异度，需要结合多个指标构建综合判别规则，并通过ROC分析评价分类性能。
